\chapter{Composite systems, density matrices and entanglement}
\label{ch:composite}

\section*{Chapter overview}

Up to this point we have described isolated systems with a single
ket. To talk about a qubit coupled to a resonator, two qubits in a
gate, or a qubit interacting with its environment, we need three
pieces of new machinery: (i) the tensor product, which combines two
Hilbert spaces into one; (ii) the density matrix, which extends the
notion of state to situations with classical uncertainty or
correlations with inaccessible degrees of freedom; and (iii) the
partial trace, which extracts the effective state of a subsystem from
that of a larger composite. With these tools we will see that
entanglement --- which has no classical analogue --- is the generic
behaviour of interacting quantum systems and the resource that powers
both quantum information processing and the standard tomography
techniques used in cQED.

% --------------------------------------------------------------
\section{Composite systems and tensor products}
\label{sec:tensor}

\begin{phenomenon}
Two independent quantum systems prepared in states $\ket{\psi_A}$
and $\ket{\psi_B}$ are jointly described by a single state in a
larger Hilbert space, the tensor product $\hilbert_A\otimes\hilbert_B$.
This much is intuitive: two non-interacting qubits, for instance,
can clearly be described by listing the state of each. The
genuinely new feature is that the tensor product is \emph{larger}
than the set of such product descriptions: it contains many states
that cannot be written as $\ket{\psi_A}\otimes\ket{\psi_B}$ for
\emph{any} choice of $\ket{\psi_A}$ and $\ket{\psi_B}$. Such states
are called \emph{entangled}; they encode correlations between the
two subsystems that have no classical analogue. The example below
shows concretely that even one of the simplest-looking two-qubit
states, $\tfrac{1}{\sqrt 2}(\ket{00}+\ket{11})$, cannot be
factorised --- any attempt to write it as a product of
single-qubit states leads to a direct contradiction.
\end{phenomenon}

\begin{postulate}[Composite systems]
\label{post:composite}
The Hilbert space of a composite system $AB$ is the tensor product
$\hilbert_{AB}=\hilbert_A\otimes\hilbert_B$
(\cref{def:tensor-product}). If $\{\ket{i}_A\}$ and $\{\ket{j}_B\}$
are orthonormal bases of $\hilbert_A,\hilbert_B$, then
$\{\ket{i}_A\otimes\ket{j}_B\}$ is an orthonormal basis of
$\hilbert_{AB}$. Operators acting only on $A$ or $B$ are extended by
the identity: $\op{O}_A\to\op{O}_A\otimes\id_B$,
$\op{O}_B\to\id_A\otimes\op{O}_B$.
\end{postulate}

\begin{example}[Two qubits and the Bell states]
\label{ex:bell-states}
For two qubits, $\hilbert_{AB}=\C^2\otimes\C^2=\C^4$ with basis
$\{\ket{00},\ket{01},\ket{10},\ket{11}\}$. The four \emph{Bell
states}
\begin{equation}
\ket{\Phi^\pm} = \tfrac{1}{\sqrt2}(\ket{00}\pm\ket{11}),
\qquad
\ket{\Psi^\pm} = \tfrac{1}{\sqrt2}(\ket{01}\pm\ket{10}),
\end{equation}
form a maximally entangled basis of $\C^4$: none of them is a
product. Let us verify this for $\ket{\Phi^+}$ by direct
calculation. Suppose, for contradiction, that
\begin{equation}
\ket{\Phi^+} \;=\;
(\alpha\ket{0}+\beta\ket{1})\otimes(\gamma\ket{0}+\delta\ket{1})
\;=\; \alpha\gamma\,\ket{00} + \alpha\delta\,\ket{01}
   + \beta\gamma\,\ket{10} + \beta\delta\,\ket{11},
\end{equation}
for some $\alpha,\beta,\gamma,\delta\in\C$. Matching coefficients
against $\ket{\Phi^+}=\tfrac{1}{\sqrt2}(\ket{00}+\ket{11})$ gives
\begin{equation}
\alpha\gamma=\tfrac{1}{\sqrt2},\quad
\beta\delta=\tfrac{1}{\sqrt2},\quad
\alpha\delta=0,\quad
\beta\gamma=0.
\end{equation}
The first two equations require $\alpha,\gamma,\beta,\delta$ all
non-zero; the last two then force $\alpha\delta=0$, a
contradiction. There are no single-qubit states whose tensor
product is $\ket{\Phi^+}$, so $\ket{\Phi^+}$ is genuinely
entangled. The same argument disposes of the other three Bell
states.
\end{example}

\paragraph{Transition.}
The four Bell states already exhibit the central oddity of composite
systems: a perfectly known joint state can correspond to two
\emph{individually} undefined subsystem states. To make sense of this,
we now generalise the notion of state itself.

% --------------------------------------------------------------
\section{Density matrices}
\label{sec:density}

\begin{phenomenon}
We rarely know the exact pure state of a real system: the
preparation procedure may produce $\ket{\psi_n}$ with classical
probability $p_n$, the system may be entangled with a bath, or the
experimenter may simply have a coarse-grained description of the
state. In all such cases the relevant object is a probability-weighted
mixture of pure states, the density matrix.
\end{phenomenon}

\paragraph{Beyond the qubit case.}
We already met the density matrix as the right description of a
single qubit when its preparation is uncertain
(\cref{def:density-tls} in \cref{sec:bloch}). The construction
makes sense in any Hilbert space, not just $\C^2$, and we restate
it here in full generality. The defining conditions
\emph{Hermitian, positive, unit-trace} are unchanged from
\cref{def:density-tls}; what is new is that $\hilbert$ is now
arbitrary --- in particular, it can be a tensor product
$\hilbert_A\otimes\hilbert_B$ of two subsystems, the case which
will let us describe a qubit's effective state when it is
entangled with another system. The single-qubit Bloch
parametrisation (\cref{stmt:bloch}) is the
$\hilbert=\C^2$ specialisation of what follows.

\begin{definition}[Density matrix, general]
\label{def:density}
A \emph{density matrix} on a Hilbert space $\hilbert$ is a linear
operator $\op{\rho}:\hilbert\to\hilbert$ such that
$\op{\rho}^\dagger=\op{\rho}$, $\op{\rho}\geq 0$, and
$\Tr\op{\rho}=1$. The set of all such operators is convex; its
extreme points are the \emph{pure states}
$\op{\rho}=\ket{\psi}\bra{\psi}$ and all other density matrices are
\emph{mixed}.
\end{definition}

For a statistical mixture of pure states $\{(\ket{\psi_n},p_n)\}$ the
density matrix is
\begin{equation}
\op{\rho} \;=\; \sum_n p_n\,\ket{\psi_n}\bra{\psi_n},
\qquad p_n\geq0,\;\sum_n p_n=1.
\end{equation}
The purity introduced for the qubit case in \cref{def:density-tls}
extends without change: $\Tr\op{\rho}^2=1$ exactly when $\op{\rho}$
is pure, and $\Tr\op{\rho}^2<1$ otherwise.

\begin{statement}[Born rule and dynamics, density-matrix form]
\label{stmt:rho-rules}
The expectation value of an observable $\op{A}$ in state $\op{\rho}$
is
\begin{equation}
\expval{\op{A}} \;=\; \Tr(\op{\rho}\,\op{A}),
\end{equation}
and unitary evolution under a Hamiltonian $\op{H}$ obeys the
Liouville--von~Neumann equation
\begin{equation}
\label{eq:liouville}
i\hbar\,\dv{t}\op{\rho} \;=\; [\op{H},\op{\rho}].
\end{equation}
\end{statement}

\begin{proof}
For a pure state $\ket{\psi}$, $\Tr(\op{\rho}\op{A})=\bra{\psi}\op{A}\ket{\psi}$
by inserting the resolution of identity; the result extends by
linearity to mixtures. For dynamics,
$\op{\rho}(t)=\op{U}\op{\rho}(0)\op{U}^\dagger$ with
$\op{U}=e^{-i\op{H}t/\hbar}$; differentiating gives
\cref{eq:liouville}.
\end{proof}

\paragraph{Connection to the Schr\"odinger equation.}
The dynamical part of \cref{stmt:rho-rules} is the Schr\"odinger
equation rewritten in density-matrix language. A pure-state
evolution under a Hamiltonian $\op{H}$ is
$\ket{\psi(t)}=\op{U}(t)\ket{\psi(0)}$ with
$\op{U}(t)=e^{-i\op{H}t/\hbar}$ (\cref{post:schrodinger}). The
corresponding density matrix therefore evolves as
\begin{equation}
\op{\rho}(t) \;=\; \ket{\psi(t)}\bra{\psi(t)}
              \;=\; \op{U}(t)\,\op{\rho}(0)\,\op{U}^\dagger(t),
\end{equation}
and differentiating both sides while using
$i\hbar\,\dot{\op{U}}=\op{H}\op{U}$ gives the Liouville--von~Neumann
equation \cref{eq:liouville}. By linearity the same equation holds
for any \emph{mixture} $\op{\rho}=\sum_n p_n\,\ket{\psi_n}\bra{\psi_n}$
of pure states evolving under the same $\op{H}$, since each
$\ket{\psi_n}\bra{\psi_n}$ obeys it term by term.

So the Liouville--von~Neumann equation is not a new dynamical
postulate: it is the Schr\"odinger equation expressed in the
language that also handles classical uncertainty in the
preparation. The two are strictly equivalent for an isolated
system. What is genuinely new comes in Part~II: when the system is
\emph{open} (coupled to an environment) we cannot describe its
dynamics by a single Hamiltonian, and the density matrix evolves
under a generally non-unitary equation --- the Lindblad master
equation --- that has no analogue at the level of a state vector,
because the state can stop being pure even if it started so.

\paragraph{Transition.}
A density matrix on $\hilbert_{AB}$ describes the joint state. The
question of what state subsystem $A$ effectively sees, when $B$ is
inaccessible, is answered by the partial trace.

% --------------------------------------------------------------
\section{The partial trace and reduced states}
\label{sec:partial-trace}

\begin{definition}[Partial trace]
\label{def:partial-trace}
Let $\op{\rho}_{AB}$ be a density matrix on
$\hilbert_A\otimes\hilbert_B$ and $\{\ket{j}_B\}$ a basis of
$\hilbert_B$. The \emph{reduced density matrix} of $A$ is
\begin{equation}
\label{eq:partial-trace}
\op{\rho}_A \;=\; \Tr_B\,\op{\rho}_{AB}
\;=\; \sum_j \bra{j}_B\,\op{\rho}_{AB}\,\ket{j}_B.
\end{equation}
$\op{\rho}_A$ is the unique operator on $\hilbert_A$ such that
$\Tr_A(\op{\rho}_A\,\op{O}_A)=\Tr(\op{\rho}_{AB}\,(\op{O}_A\otimes\id_B))$
for every observable $\op{O}_A$ on $A$.
\end{definition}

\begin{statement}[Reduction of an entangled pure state]
\label{stmt:reduced-pure}
Let $\ket{\Psi}_{AB}$ be a pure state. Its reduced state
$\op{\rho}_A=\Tr_B\,\ket{\Psi}\bra{\Psi}$ is mixed if and only if
$\ket{\Psi}$ is entangled. In particular,
\begin{equation}
\Tr\op{\rho}_A^2 = 1 \;\;\Longleftrightarrow\;\;
\ket{\Psi}_{AB} = \ket{\psi}_A\otimes\ket{\phi}_B
\text{ for some }\ket{\psi},\ket{\phi}.
\end{equation}
\end{statement}

\begin{proof}[Sketch via Schmidt decomposition]
Any bipartite pure state admits the Schmidt decomposition
$\ket{\Psi}=\sum_k \lambda_k\,\ket{u_k}_A\otimes\ket{v_k}_B$ with
$\lambda_k\geq0$ and $\sum_k\lambda_k^2=1$. Tracing over $B$ yields
$\op{\rho}_A=\sum_k\lambda_k^2\,\ket{u_k}\bra{u_k}$, which is pure
($\lambda_k=\delta_{k1}$) precisely when only one Schmidt coefficient
is non-zero, i.e.\ when $\ket{\Psi}$ factorises.
\end{proof}

\begin{example}[Bell state has maximally mixed marginals]
For $\ket{\Phi^+}=(\ket{00}+\ket{11})/\sqrt2$ a direct calculation
gives
$\op{\rho}_A = \tfrac12(\ket{0}\bra{0}+\ket{1}\bra{1}) = \tfrac12\id$.
Each qubit individually is in the maximally mixed state $\id/2$ ---
its Bloch vector is the origin of the Bloch sphere --- even though
the joint state is perfectly known.
\end{example}

\begin{remark}
The same construction underlies the description of an open system
coupled to a bath: the system's effective state is the partial trace
of the system--bath density matrix over the bath. The Lindblad
master equation of \cref{ch:open-systems} is the time-local equation
this reduced state obeys under the Born--Markov approximation.
\end{remark}

\paragraph{Transition.}
Entanglement, having appeared as the obstruction to writing
$\ket{\Psi}$ as a product, deserves a section of its own --- both as
a phenomenon and as a quantitative resource.

% --------------------------------------------------------------
\section{Entanglement}
\label{sec:entanglement}

\begin{definition}[Separable and entangled states]
\label{def:entanglement}
A pure bipartite state $\ket{\Psi}_{AB}$ is \emph{separable} if it
factorises as $\ket{\psi}_A\otimes\ket{\phi}_B$, otherwise it is
\emph{entangled}. A mixed state $\op{\rho}_{AB}$ is separable if it
admits a decomposition
$\op{\rho}_{AB}=\sum_n p_n\,\op{\rho}_A^{(n)}\otimes\op{\rho}_B^{(n)}$
with $p_n\geq0$, $\sum_n p_n=1$, and entangled otherwise.
\end{definition}

\begin{statement}[Schmidt rank as an entanglement witness]
\label{stmt:schmidt}
The Schmidt rank of a bipartite pure state $\ket{\Psi}_{AB}$ ---
i.e.\ the number of non-zero Schmidt coefficients $\lambda_k$ in
$\ket{\Psi}=\sum_k\lambda_k\ket{u_k}_A\ket{v_k}_B$ --- equals one if
and only if $\ket{\Psi}$ is separable. Larger Schmidt rank means
more entanglement; the entropy $S(\op{\rho}_A)=-\sum_k\lambda_k^2\log\lambda_k^2$
of the reduced state is the canonical quantitative measure.
\end{statement}

\begin{example}[Two-qubit entanglement: one ``ebit'' in $\ket{\Phi^+}$]
We compute the entanglement quantitatively for the Bell state of
\cref{ex:bell-states}. By inspection, $\ket{\Phi^+}$ is already in
Schmidt form,
\begin{equation}
\label{eq:phiplus-schmidt}
\ket{\Phi^+} \;=\; \tfrac{1}{\sqrt2}\,\ket{0}_A\!\otimes\!\ket{0}_B
                 + \tfrac{1}{\sqrt2}\,\ket{1}_A\!\otimes\!\ket{1}_B,
\end{equation}
with two non-zero Schmidt coefficients
$\lambda_1=\lambda_2=\tfrac{1}{\sqrt2}$. Because the two
$\lambda_k$ are equal, the bipartite state is \emph{maximally
entangled} on $\C^2\otimes\C^2$: any other choice of unit-norm
$\lambda_k$ with $\lambda_1^2+\lambda_2^2=1$ has at least one
coefficient strictly larger than $1/\sqrt 2$ and is more
``product-like'' in the sense that subsystem $A$ retains more
information about its own state.

\paragraph{The reduced state.}
Applying the partial-trace formula \cref{eq:partial-trace} to
\cref{eq:phiplus-schmidt} (or specialising the proof of
\cref{stmt:reduced-pure} to this case) gives
\begin{equation}
\op{\rho}_A \;=\; \Tr_B\,\ket{\Phi^+}\bra{\Phi^+}
              \;=\; \lambda_1^2\,\ket{0}\bra{0}
                  + \lambda_2^2\,\ket{1}\bra{1}
              \;=\; \tfrac12\,\id,
\end{equation}
i.e.\ qubit $A$ on its own is in the maximally mixed state ---
Bloch vector at the origin. Even though the joint $\ket{\Phi^+}$
is perfectly known, an observer with access only to $A$ has zero
predictive information about $A$'s observables: this is the
content of \cref{stmt:reduced-pure} made concrete.

\paragraph{The entanglement entropy.}
Plug $\lambda_1=\lambda_2=1/\sqrt 2$ into the entropy formula of
\cref{stmt:schmidt}:
\begin{equation}
S(\op{\rho}_A)
\;=\; -\sum_k \lambda_k^2\,\log_2 \lambda_k^2
\;=\; -2\cdot\tfrac12\,\log_2\tfrac12
\;=\; \log_2 2 \;=\; 1\;\text{bit}.
\end{equation}
This is the maximum value $S(\op{\rho}_A)$ can take on
$\C^2\otimes\C^2$ (it is bounded by $\log_2\dim\hilbert_A=1$ for a
qubit), confirming that $\ket{\Phi^+}$ is maximally entangled.

\paragraph{The unit ``ebit''.}
The result $S(\op{\rho}_A)=1$ bit motivates a unit of its own. An
\emph{ebit} (entanglement bit) is the conventional unit of
bipartite pure-state entanglement, defined as the amount of
entanglement contained in one Bell pair: a state with
$S(\op{\rho}_A)=n$ bits carries $n$ ebits. The unit is more than
notational --- under local operations and classical communication
(LOCC), bipartite pure entanglement is, asymptotically, freely
convertible into ebits and back, so any pure-state entanglement
resource can be quantified in this single unit. Quantum
teleportation of one qubit consumes exactly one ebit; the
production of $\ket{\Phi^+}$ by an idealised CNOT gate acting on
$\ket{+}_A\otimes\ket{0}_B
= \tfrac{1}{\sqrt 2}(\ket{0}+\ket{1})_A\!\otimes\!\ket{0}_B$
\emph{produces} one ebit, since
\begin{equation}
\text{CNOT}\,\bigl(\tfrac{1}{\sqrt 2}(\ket{0}+\ket{1})_A\!\otimes\!\ket{0}_B\bigr)
\;=\; \tfrac{1}{\sqrt 2}(\ket{00}+\ket{11})
\;=\; \ket{\Phi^+}.
\end{equation}
This is also the motivation for calling two-qubit entangling gates
``the ebit-generating part of any quantum-computational primitive''
(\cref{app:entanglement-cqed}).
\end{example}

\begin{application}[Why entanglement matters in cQED]
\label{app:entanglement-cqed}
Two-qubit gates (\cref{sec:two-qubit-gates}) are characterised by
their ability to entangle: a gate that cannot create entanglement is
locally equivalent to a product of single-qubit gates and provides
no quantum-computational power. Conversely, entanglement with the
environment is the cause of decoherence
(\cref{ch:open-systems}): tracing out the environment in an
entangled system--bath state produces a mixed reduced state for the
qubit.
\end{application}

% --------------------------------------------------------------
\section*{Chapter summary and outlook}

We have moved from a single isolated system described by a ket to
a much richer setting: composite systems live in tensor products,
mixtures and reductions are described by density matrices, and the
partial trace makes the relationship between joint and marginal
states explicit. Entanglement emerges as the generic property of
non-product joint states, and shows up in cQED both as the resource
exploited by gates and as the mechanism behind decoherence.

The final chapter of this part deals with one further question that
the postulates do not, by themselves, answer: how does a system
respond to a small or near-resonant time-dependent drive? The
answer assembles two of the workhorses of quantum optics ---
time-dependent perturbation theory (and Fermi's golden rule) and the
rotating-wave approximation --- both of which we will use repeatedly
in Parts~III and IV.
